% !TeX root = ../数模通用模板.tex
\section{符号说明}
公共符号见表\ref{tab:global-symbols}，回归系数及各问局部辅助变量在使用处定义。$E$表示内部储电量，$e$表示紧急购电量。
\begin{table}[H]\centering
\caption{主要符号及单位}\label{tab:global-symbols}
\renewcommand{\arraystretch}{1.1}
\begin{tabularx}{\textwidth}{p{3.1cm}Xc}\toprule
符号 & 含义 & 单位\\\midrule
$d,t$ & 日期与日内槽索引，$t=1,\ldots,144$ & --\\
$\Delta t$ & 槽长，取$1/6$ & h\\
$\tau$ & 预报发布小时，取0、6、12、18 & h\\
$L_{d,t},V_{d,t}$ & 区间平均负载、光伏功率 & kW\\
$p_{d,t},\widehat p_{d,t}$ & 实际基础电价与预测电价 & 元/kWh\\
$g_{d,t}$ & 午夜原计划购电量 & kWh\\
$q_{d,t}$ & 该槽最后有效版本的常规购电总量 & kWh\\
$a^+_{d,t},a^-_{d,t}$ & 相对午夜原计划的最终净上调、净下调量 & kWh\\
$e_{d,t}$ & 实际紧急购电量 & kWh\\
$c_{d,t},d_{d,t}$ & 交流侧充电、放电量 & kWh\\
$P^c_{d,t},P^d_{d,t}$ & 交流侧充电、放电功率 & kW\\
$E_{d,t}$ & 第$t$槽开始时的内部储电量 & kWh\\
$E_{\min},E_{\max}$ & 内部储电量下、上界 & kWh\\
$P_{\max}$ & 单方向最大功率 & kW\\
$\eta_c,\eta_d$ & 单向充、放电效率，均为$\sqrt{0.9}$ & --\\
$w_{d,t}$ & 总未利用电量 & kWh\\
$\mathcal I_{d,\tau}$ & 购电发布时刻可用信息集 & --\\
$s,\pi_s$ & 情景索引与权重 & --\\
$C$ & 实际购电费用 & 元\\
$J$ & 各问规划代理目标 & 见各问定义\\\bottomrule
\end{tabularx}\end{table}
功率与交流侧电量按$c=P^c\Delta t$、$d=P^d\Delta t$转换。以$n=(L-V)\Delta t$表示净需求电量，实际执行满足
\begin{equation}\label{eq:common-balance}
 q+d+e-c-w=n,\qquad E_{t+1}=E_t+\eta_c c_t-d_t/\eta_d.
\end{equation}
普通购电量$q$不含紧急购电，外网总购电量为$q+e$。问题二和4-2中$q=g$；问题三和4-3中，未执行区间采用最后一次有效更新，令$a^+=(q-g)^+$、$a^-=(g-q)^+$，实际费用统一为
\begin{equation}\label{eq:common-bill}
 C=\sum_{d,t}p_{d,t}\bigl(g_{d,t}+1.5a^+_{d,t}-0.5a^-_{d,t}+5e_{d,t}\bigr).
\end{equation}
工作簿的“调整购电量”填写$q$，不是$a^+-a^-$，不再与$g$重复相加。吞吐、启动、功率变化惩罚与日末库存估值只进入$J$，不属于式\eqref{eq:common-bill}的真实账单。文中费用保留2位小数、电量保留4位小数，差额和比例均由未舍入结果计算。
